\section{Locking small overlaps across a large overlap}
\label{static-sec:locking}

When two configurations have a positive overlap bounded away from zero,
we will show that their small overlaps with a third configuration are
close, including their signs. This is the fixed-scale estimate used by
the crossing construction. Spin reversal will then give the
absolute-overlap version needed in Appendix~\ref{fixed-sec:crossings}.

Throughout the section, fix $\beta>1$ and a minimizer $\mu$ of the scalar
functional. We use the notation of Section~\ref{static-sec:scalar} and
write $\pi=\pi_{n,\beta}^J$. The scalar inputs are
\begin{equation}
\label{static-eq:locking-scalar-identities}
\begin{gathered}
0\leq\chi\leq1,\qquad |u|\leq1,\qquad
|\chi(s,z)-\chi(s,z')|\leq C_\beta|z-z'|,\\
\E\chi(s,X_s)=A(s)+\alpha(s)\bigl(s-\Gamma(s)\bigr),\qquad
\E\bigl[\chi(s,X_s)+\alpha(s)u(s,X_s)^2\bigr]=d_*(s),\\
\beta A(0)\leq1,\qquad A(h)<A(0)\quad(0<h\leq1),\qquad
d_*(s)\longrightarrow A(0)\quad(s\downarrow0).
\end{gathered}
\end{equation}
We also use the scalar martingale evolution from
Lemma~\ref{static-lem:scalar-regularity} and the bounds
$\alpha(s)>0$ for $s>0$ and $0\le d_*(s)\le1$ from
Proposition~\ref{static-prop:scalar-minimizer}.

For a nonempty Gram constraint $D$, write $Z_{d,n}(D)$ for the
unnormalized partition sum in~\eqref{static-eq:constrained-pressure}.
Thus $F_{d,n}(D)=n^{-1}\E\log Z_{d,n}(D)$, and division by
$Z_{n,\beta}^d$ gives the probability of $D$ under $d$ independent
Gibbs draws.

\begin{lemma}[Signed fixed-scale locking]\label{lem:fixed-locking}
Let $e_*=1/100$. For every fixed $h\in(0,1)$, there is $b>0$ such that,
for all sufficiently large $n$,
\begin{equation}\label{eq:fixed-locking}
 \E\pi^{\otimes3}\!\left\{
 \begin{gathered}
 R(x,y)\ge h,\qquad |R(v,x)|,|R(v,y)|\le b,\\
 |R(v,x)-R(v,y)|>2n^{-e_*}
 \end{gathered}\right\}\le e^{-n^{9/10}}.
\end{equation}
The same conclusion holds after decreasing $b$.
\end{lemma}

We will prove a pressure deficit for each forbidden overlap matrix and
then sum probability bounds. The conversion below will also apply to
the shrinking gaps in Section~\ref{sec:narrow}.

\begin{lemma}[From a pressure gap to a probability bound]
\label{lem:pressure-prob}
Fix a replica number $d$. If a feasible overlap matrix $D$ satisfies
\[
 F_{d,n}(D)\le dP_\beta-\eta,
 \qquad \eta\ge2dC_\beta n^{-1/24},
\]
where $C_\beta$ is the constant in
Theorem~\ref{static-thm:finite-size}, then, for all sufficiently large $n$,
\begin{equation}
\label{static-eq:pressure-gap-probability}
 \E\frac{Z_{d,n}(D)}{Z_{n,\beta}^d}
 \le \exp(-n\eta/4)+\exp(-c_{\beta,d}n\eta^2).
\end{equation}
The constants are uniform over $D$ and over $n$-dependent gaps satisfying
the displayed lower bound.
\end{lemma}
\begin{proof}
For fixed $d$, the derivative of $\log Z_{d,n}(D)$ with respect to any
$J_{ij}$ has absolute value at most $\beta d/\sqrt n$. Both this function
and $d\log Z_{n,\beta}$ therefore have Gaussian Lipschitz constant at most
$\beta d\sqrt{(n-1)/2}$. Their difference
\[
 T_D:=\log Z_{d,n}(D)-d\log Z_{n,\beta}
\]
has Lipschitz constant at most $C_{\beta,d}\sqrt n$ and satisfies
$T_D\le0$, since $e^{T_D}$ is a Gibbs probability.
Theorem~\ref{static-thm:finite-size} and the pressure-gap hypothesis give
$\E T_D\le-n\eta/2$. Gaussian concentration bounds
$\Prob\{T_D>-n\eta/4\}$ by $\exp(-c_{\beta,d}n\eta^2)$.
On its complement $e^{T_D}\le e^{-n\eta/4}$, and on the event itself
$e^{T_D}\le1$. Summing proves the claim. This calculation uses the
original off-diagonal disorder, so no auxiliary diagonal Gaussian
appears in the probability bound.
\end{proof}

Only $d=2,3$ will be needed. There are at most
$(n+1)^{d(d-1)/2}$ feasible overlap matrices, so a union over all
constraints costs only a polynomial factor.

\subsection{A pair constraint}

The three-replica comparison will use a large overlap $q$ for which
$\Gamma(q)$ is close to $q$. We first show that the remaining pair
overlaps already have an exponentially small Gibbs probability.

\begin{lemma}[Pair pressure bound]
\label{static-lem:pair-bound}
For every nonempty pair constraint with overlap $q\in[0,1]$,
\begin{equation}
\label{static-eq:pair-pressure}
 F_{2,n}\!\left(\begin{smallmatrix}1&q\\q&1\end{smallmatrix}\right)
 \leq 2P_\beta-\frac12\bigl(\Gamma(q)-q\bigr)^2.
\end{equation}
Consequently, for every fixed $\tau>0$, the averaged Gibbs probability
of pairs satisfying
$\bigl|\Gamma(|R(x,y)|)-|R(x,y)|\bigr|>\tau$
is at most $C\exp(-cn)$ for all large $n$.
\end{lemma}

\begin{proof}
Use Proposition~\ref{static-prop:matrix-bound} with two coordinates
sharing their covariance increments up to scalar time $q$ and receiving
independent increments thereafter. Give the common increments masses
$\alpha(s)/2$ and the independent increments masses $\alpha(s)$.
The endpoint is the prescribed Gram matrix.

At zero terminal multiplier this path has recursive value
$2\Phi(0,0)$. Indeed, integrating the independent portion first leaves
$\Phi(q,z_x)+\Phi(q,z_y)$. On the earlier common portion $z_x=z_y$;
integration of twice a scalar value at mass $\alpha/2$ is twice the
scalar integration at mass $\alpha$. The penalty has the same factor
two: before the split the squared norm is $4s^2$ with mass $\alpha/2$,
and after the split its varying part is $2s^2$ with mass $\alpha$.

Now add $L s_xs_y$ to the terminal exponent and denote the recursive
value by $\Psi(L)$. In the matrix convention this sets
$\lambda_{xy}=\lambda_{yx}=L/2$, so $\lambda:D=Lq$. Differentiating at
$L=0$ transports the terminal observable $s_xs_y$ backwards by
conditional expectation under the unperturbed tilted path. Given the
common field at $q$, the two later branches are independent and each
has conditional spin mean $u(q,X_q)$. Hence
\[
 \Psi'(0)=\E u(q,X_q)^2=\Gamma(q).
\]
We also need the bound $\Psi''(L)\leq1$ for every $L$. For a terminal
observable $O\in[-1,1]$, the terminal Hessian is its thermal variance.
Backward integration at mass $m\in[0,1]$ averages the existing Hessian
and adds $m$ times the conditional variance of the gradient. Iterating
this formula and using the total-variance decomposition bounds the
result by $\Var(O)\leq1$ under the tilted path law. An outer quenched
expectation preserves the bound. Taylor's inequality therefore gives
\[
 F_{2,n}(q)\leq2P_\beta+L\bigl(\Gamma(q)-q\bigr)+\frac{L^2}{2}.
\]
The choice $L=q-\Gamma(q)$ yields~\eqref{static-eq:pair-pressure}.

These computations are first made with step approximations of $\alpha$
whose segment endpoints include $q$. Lemma~\ref{static-lem:scalar-regularity}
gives convergence of the scalar values, derivatives, and tilted
transitions, so the inequality passes to the minimizing $\alpha$.

Finally, reversing one spin configuration preserves its Gibbs weight
and changes the sign of its overlap. The same pressure estimate thus
applies with the absolute overlap in place of $q$. There are at most
$n+1$ pair overlaps; summing~\eqref{static-eq:pressure-gap-probability}
with the fixed deficit $\tau^2/2$ proves the probability assertion.
\end{proof}

\subsection{A three-coordinate path with an asymmetric perturbation}
\label{static-subsec:locking-base}

Fix $q\in[h,1]$ and a nonempty constraint
\begin{equation}
\label{static-eq:locking-gram}
 D=\begin{pmatrix}
 1&r+\delta&r-\delta\\
 r+\delta&1&q\\
 r-\delta&q&1
 \end{pmatrix},\qquad r\geq0.
\end{equation}
We will choose upper bounds for $r$ and $|\delta|$ depending only on
$\beta,h$. Write $D_0$ for the same matrix with $\delta=0$. Starting
from a hierarchy ending at $D_0$, we will add opposite field increments
to the $x$ and $y$ coordinates. This realizes the desired first-order
change in their overlaps with $v$. Keeping the increments positive
semidefinite also adds a second-order covariance term. The free-endpoint
matrix bound charges only its squared size, so that correction has a
fourth-order cost.

The reason to perturb an early increment is the strict inequality
$\beta^2A(h)A(0)<1$, supplied by
\eqref{static-eq:locking-scalar-identities}. At a small branching time
$r$, the two scalar branches are nearly independent. Their
susceptibility product is then bounded above by $A(h)A(0)$, up to
small errors.
The Gaussian variance part of the perturbation will cancel the
corresponding penalty up to the pair deviation $q-\Gamma(q)$; this
strict product inequality controls what remains.

The base hierarchy itself is valid for every $0\le r\le q\le1$;
only the later perturbation will require $r$ to be small. We use the
three-replica construction in Talagrand's discussion following
\cite[Proposition~4.4]{Talagrand07}, with masses divided by the size
of each common block. Its three stages, in local scalar time $s$, are:
\begin{enumerate}[label=\textup{(\roman*)}]
\item From $0$ to $r$, all three coordinates share each increment,
      integrated at mass $\alpha(s)/3$.
\item From $r$ to $q$, the coordinate $v$ follows an independent branch
      at mass $\alpha(s)$, while $x,y$ share a branch at mass
      $\alpha(s)/2$.
\item From $q$ to $1$, all three coordinates have independent increments
      at mass $\alpha(s)$.
\end{enumerate}
In the middle stage each branch has its own clock. If these clocks are
$a$ for $v$ and $b$ for the common $x,y$ branch, the current covariance is
\[
 Q(a,b)=\begin{pmatrix}a&r&r\\r&b&b\\r&b&b\end{pmatrix},
 \qquad r\le a,b\le q.
\]
Thus one branch can advance while the other is held fixed. For a step
function $\alpha$, merge the two finite lists of increments in nondecreasing
order of mass, retaining the order within each list. Each list is
nondecreasing, so such a merge exists; ties may be ordered in either
way consistent with the two local orders. The first-stage masses do
not exceed either middle starting mass, and the middle masses do not
exceed the starting mass of the last stage. Thus the entire schedule
is a valid nondecreasing hierarchy. Insert a dummy initial covariance
$Q_0=0$ to ensure that, even when $\alpha(0)>0$, the first genuine field
increment has its prescribed positive mass and the quenched initial
level has zero covariance.

The endpoint is $D_0$, and the recursive value minus the penalty equals
$3\cP_\beta(\mu)$. To check both assertions about the value, observe
that a common block of size $\ell$ has scalar value $\ell\Phi$.
Its mass $\alpha/\ell$ gives the ordinary scalar tilted transition,
and its squared-norm differential $\ell^2\dd(s^2)$ contributes
$\ell$ times the scalar penalty after multiplication by that mass.
Independent branches have additive values. When one branch is
integrated, the value of a frozen branch is independent of that
increment and passes unchanged through the logarithm. Applying these
facts at each split proves the claimed value and penalty for every
allowed interleaving. It also identifies the base tilted law. Denote the $v$ path and the
shared $x,y$ path by $X^v$ and $X^x$, and their common field at the
first split by $X_r$. Write $K_{r,z}$ for the scalar path law started at $(r,z)$ and run to time $q$.
If $\mathcal X_0$ denotes the common prefix through $r$, then
\begin{equation}\label{static-eq:locking-product-law}
 \mathcal L(X^v_{[r,q]},X^x_{[r,q]}\mid\mathcal X_0)
   =K_{r,X_r}\otimes K_{r,X_r}.
\end{equation}
Indeed, the frozen branch cancels from each normalized kernel, while
the factor three or two in the common-block value cancels the divided
mass. The two future branches are therefore independent conditional on
$X_r$, for every allowed merge. We will use this same base hierarchy
in Section~\ref{sec:narrow}.

For the fixed-scale perturbation, now fix $0<s_1<s_2<h$ and suppose
$r<s_1$. Define
\[
 k(s)=\frac{\1_{[s_1,s_2]}(s)}{s_2-s_1},\qquad
 g(s)=\int_0^s k(t)\dd t,\qquad
 G(s)=\int_0^s k(t)^2\dd t,\qquad
 K=G(s_2)=\frac1{s_2-s_1}.
\]
Let $L_{vx}=L_{xv}=1$, $L_{vy}=L_{yv}=-1$, with all other entries
of $L$ zero. Let $L_2$ vanish outside its $(x,y)$ block and give that
block the value
$\left(\begin{smallmatrix}1&-1\\-1&1\end{smallmatrix}\right)$.
Then
\[
 |L|^2=|L_2|^2=4,\qquad L:L_2=0.
\]
During each $v$-increment with local clock in $[s_1,s_2]$, add
$\delta k$ times that increment to $x$ and its negative to $y$.
At the covariance level, this replaces $\dd s\,e_ve_v^{\mathsf T}$ by
\[
 \dd s\,
 \bigl(e_v+\delta k(e_x-e_y)\bigr)
 \bigl(e_v+\delta k(e_x-e_y)\bigr)^{\mathsf T}.
\]
This rank-one form is positive semidefinite. The accumulated change
is $\delta gL+\delta^2GL_2$ until the modification ends, and is then
held fixed. The new endpoint is consequently
\[
 \widetilde D=D_0+\delta L+\delta^2KL_2
              =D+\delta^2KL_2.
\]
Apply Proposition~\ref{static-prop:matrix-bound} with zero terminal
multiplier and this endpoint. Its mismatch contribution is
\begin{equation}
\label{static-eq:locking-mismatch}
 \frac{\beta^2}{4}|D-\widetilde D|^2
 =\beta^2\delta^4K^2.
\end{equation}

We next isolate the beneficial increase in the penalty. Along the
base path, $Q:L=0$. Also $Q:L_2=0$ until the split at $q$, by which
time $g=1$ and $G=K$. In the final stage, $Q:L_2=2(s-q)$. Thus
\[
 |\widetilde Q|^2-|Q|^2
 =4\delta^2g^2+2\delta^2G(Q:L_2)+4\delta^4G^2.
\]
The last term has nonnegative increments. The first changes only
during the $v$-increments, whose masses are $\alpha(s)$; the middle
term changes only in the independent final stage. The penalty
increase is therefore at least
\begin{equation}
\label{static-eq:locking-penalty}
 \beta^2\delta^2\left(
     \int_0^1\alpha(s)\dd(g(s)^2)+K A(q)\right).
\end{equation}

\subsection{The recursive value under the perturbation}

To compare this penalty gain with the recursive value, first integrate
out the unchanged independent final stage. Let $z$ be the common base
field of $x,y$ at the end of the middle stage, and let $Y$ be the sum
of the raw $v$-increments weighted by $k$. The boundary value changes by
\begin{equation}
\label{static-eq:locking-boundary-increment}
 \Delta=\Phi(q,z+\delta Y)+\Phi(q,z-\delta Y)-2\Phi(q,z).
\end{equation}
Convexity, the bound $\chi\leq1$, and its spatial Lipschitz bound yield
\begin{equation}
\label{static-eq:locking-taylor}
 0\leq\Delta\leq\delta^2Y^2,
 \qquad
 \bigl|\Delta-\delta^2\chi(q,z)Y^2\bigr|
 \leq C_\beta|\delta Y|^3.
\end{equation}

The remaining backward integrations are nonlinear, so a bound on the
boundary change must be propagated through their normalized kernels.
If the future value changes by $U$, an integration at positive mass
$m$ changes by
\[
 \frac1m\log\E_{\rm base}
              [e^{mU}\mid\text{current prefix}],
\]
where the expectation uses the unperturbed tilted kernel. At zero
mass, the change is its conditional expectation of $U$. For
$0<m\leq1$, Jensen's inequality and the power-mean inequality put
the displayed expression between
$\E_{\rm base}[U\mid\text{prefix}]$ and
$\log\E_{\rm base}[e^U\mid\text{prefix}]$.

Here is the resulting induction explicitly. Let $K_i$ be the
normalized base kernel at level $i$, and $U_i$ the difference of
recursive values at that prefix. Taking the ratio of the two defining
integrals gives $U_i=m_i^{-1}\log K_i e^{m_iU_{i+1}}$ for $m_i>0$;
this identity uses the base kernel only. Put
\[
 A_i=\E_{\rm base}[\Delta\mid\text{prefix }i],\qquad
 B_i=\log\E_{\rm base}[e^\Delta\mid\text{prefix }i].
\]
At the boundary these equal $U_i=\Delta$. If
$A_{i+1}\leq U_{i+1}\leq B_{i+1}$, the kernel inequalities and tower
property give
\[
 U_i\geq K_iU_{i+1}\geq A_i,\qquad
 U_i\leq\log K_ie^{U_{i+1}}
       \leq\log K_ie^{B_{i+1}}=B_i.
\]
Jensen's inequality gives the same conclusions at zero mass. Applying
this induction through the hierarchy proves
\begin{equation}
\label{static-eq:locking-recursive-sandwich}
 \E_{\rm base}\Delta
 \ \leq\ \widetilde\Psi-\Psi
 \ \leq\ \log\E_{\rm base}e^\Delta.
\end{equation}
Jensen's inequality also preserves these bounds under an outer
quenched expectation.

Under the base tilted law, take scalar diffusions $X^v,X^x$ that agree
through time $r$ and have independent Brownian noises thereafter.
The identified path law gives
\[
 z=X_q^x,\qquad
 Y=\int_{s_1}^{s_2} k(s)\dd X_s^v,
 \qquad
 \dd X_s^v=\beta\dd B_s^v
          +\beta^2\alpha(s)u(s,X_s^v)\dd s.
\]
Thus $Y$ is a centered Gaussian increment of variance $\beta^2K$
plus a possibly dependent random term of absolute value at most
$\beta^2$. In particular, there are $c,C>0$, depending only on
$\beta,s_1,s_2$, such that
\begin{equation}
\label{static-eq:locking-Y-moment}
 \E_{\rm base}e^{cY^2}\leq C.
\end{equation}
This estimate is uniform in $q\geq h$, $r<s_1$, and the step
approximations of $\alpha$. For small enough $|\delta|$,
\eqref{static-eq:locking-taylor} and this exponential moment imply
\[
 \E\bigl(e^\Delta-1-\Delta\bigr)
 \leq\tfrac12\E\bigl[\Delta^2e^\Delta\bigr]
 \leq C\delta^4.
\]
Combining with~\eqref{static-eq:locking-recursive-sandwich} bounds the
recursive change by
\begin{equation}
\label{static-eq:locking-recursive-expansion}
 \widetilde\Psi-\Psi
 \leq\delta^2\E\bigl[\chi(q,X_q^x)Y^2\bigr]+C|\delta|^3.
\end{equation}

At this point all matrix paths can still be taken finite. Include
$r,q,s_1,s_2$ in their segment endpoints and choose nondecreasing step
approximations converging to $\alpha$ in $L^1([0,1])$.
Lemma~\ref{static-lem:scalar-regularity} gives convergence of the scalar
values, derivatives, and diffusions. Couple the diffusions with the
same Brownian noises, keeping the noises independent after $r$.
The uniform exponential moment~\eqref{static-eq:locking-Y-moment}
then justifies passage to the limit in the expectation in
\eqref{static-eq:locking-recursive-expansion}. The penalty integrals
converge as well: $\dd(g^2)/\dd s$ is bounded and supported on
$[s_1,s_2]$, and $A(q)$ is an integral of $\alpha$.
We have therefore applied the matrix comparison only to finite
hierarchies and passed their upper bounds to the limit. No infinite
hierarchy has been differentiated. In particular, convergence at a
possible atom $q$ is unnecessary; only $L^1$ convergence of the
masses is used. The minimizing identities will be applied after this
limit.

Subtracting~\eqref{static-eq:locking-penalty} from
\eqref{static-eq:locking-recursive-expansion} and including the
fourth-order cost~\eqref{static-eq:locking-mismatch} gives
\begin{equation}
\label{static-eq:locking-three-pressure-expansion}
 \begin{split}
 F_{3,n}(D)\leq 3P_\beta+\delta^2\mathcal B(q,r)+C|\delta|^3,\\
 \mathcal B(q,r)=
 \E\bigl[\chi(q,X_q^x)Y^2\bigr]
 -\beta^2KA(q)-\beta^2\int_0^1\alpha(s)\dd(g(s)^2).
 \end{split}
\end{equation}
The task is now to choose the perturbation interval so that this
quadratic coefficient is strictly negative.

\subsection{Choosing an interval with a strict quadratic gain}

The strict inequality $A(h)<A(0)$ will supply the gain. To expose it,
define
\[
 V_t=\chi(t,X_t^v)+\alpha(t)u(t,X_t^v)^2,
 \qquad Y_s=\int_{s_1}^{s}k(t)\dd X_t^v
 \quad(s\in[s_1,s_2]).
\]
The martingale identity
$\dd u(s,X_s^v)=\beta\chi(s,X_s^v)\dd B_s^v$ and It\^o's formula,
conditional on the common starting field $X_r$, give for almost every
$s$
\[
 \frac{\dd}{\dd s}\E[Y_su(s,X_s^v)\mid X_r]
   =\beta^2 k(s)\E[V_s\mid X_r].
\]
Since $Y_{s_1}=0$, apply It\^o's formula once more to $Y_s^2$ and use
conditional independence of the $v$ and $x$ branches. This yields
\begin{equation}
\label{static-eq:locking-Ito}
 \begin{split}
 \E\bigl[\chi(q,X_q^x)Y^2\bigr]
 ={}&\beta^2K\E\chi(q,X_q)\\
 &+2\beta^4\int_{s_1}^{s_2}\alpha(s)k(s)
      \int_{s_1}^{s}k(t)
       \E\bigl[\chi(q,X_q^x)V_t\bigr]\dd t\dd s.
 \end{split}
\end{equation}
All factors except $Y$ are bounded, and its required moments follow
from~\eqref{static-eq:locking-Y-moment}. The conditioning and integrations
are therefore justified.

Only the common field at the small time $r$ prevents the product
expectation in this formula from factoring. Its effect is bounded by
\begin{equation}
\label{static-eq:locking-decoupling}
 \left|\E\bigl[\chi(q,X_q^x)V_t\bigr]
       -\E\chi(q,X_q)\,d_*(t)\right|
 \leq C_\beta\sqrt r,
 \qquad s_1\leq t\leq s_2,\quad h\leq q\leq1.
\end{equation}
To prove this estimate, restart the $x$ branch at an independent copy
$X_r'$ of $X_r$, independent also of the whole $v$ branch, and retain
its post-$r$ Brownian noise. The drift $\beta^2\alpha u$ is uniformly
Lipschitz in space, so the two $x$ solutions at time $q$ differ by at
most $e^{\beta^2}|X_r-X_r'|$. Use the spatial Lipschitz bound for
$\chi$, the bound $0\leq V_t\leq2$, and
$\E|X_r-X_r'|\leq C_\beta\sqrt r$, the latter following from the
bounded drift. The replacement has the same marginal law and is
independent of $V_t$; its expectation therefore factors, with
$\E V_t=d_*(t)$ by~\eqref{static-eq:locking-scalar-identities}.

The two inequalities $\beta A(0)\leq1$ and $A(h)<A(0)$ imply
\[
 \beta^2 A(h)A(0)<1.
\]
Choose $s_2\in(0,h)$ and $\gamma>0$ small enough that
\begin{equation}
\label{static-eq:locking-margin}
 \beta^2A(h)\sup_{0<t\leq s_2}d_*(t)\leq1-4\gamma,
\end{equation}
and set $s_1=s_2/2$. The functions $k,g$ and the constant $K$ are now
fixed. Also set
\[
 W=2\beta^2\int_{s_1}^{s_2}\alpha(s)k(s)
                    \int_{s_1}^{s}k(t)\dd t\dd s>0.
\]
The strict positivity follows from $\alpha(s)>0$ for every $s>0$.

Let $\bar\chi(q)=\E\chi(q,X_q)$. For $q\geq h$ with
$|q-\Gamma(q)|\leq\tau$, the all-time scalar identity gives
\[
 \bar\chi(q)=A(q)+\alpha(q)(q-\Gamma(q))\leq A(h)+\tau.
\]
First choose $\tau>0$ small enough and then
$r_0\in(0,s_1)$ small enough. Equations
\eqref{static-eq:locking-decoupling} and
\eqref{static-eq:locking-margin} give
\begin{equation}
\label{static-eq:locking-negative-factor}
 \beta^2\E\bigl[\chi(q,X_q^x)V_t\bigr]-1\leq-2\gamma
\end{equation}
uniformly for $q\in[h,1]$ with $|q-\Gamma(q)|\leq\tau$,
$0\leq r\leq r_0$, and $t\in[s_1,s_2]$. More explicitly, since
$0\leq d_*(t)\leq1$, the left side is bounded above by
$-4\gamma+\beta^2\tau+C_\beta\sqrt r$.

We can now compare the entire quadratic term with this strict margin.
Using $\dd(g^2)=2gk\dd s$, substitute~\eqref{static-eq:locking-Ito}
into~\eqref{static-eq:locking-three-pressure-expansion}. The result is
the exact identity
\begin{equation}
\label{static-eq:locking-bracket}
 \begin{split}
 \mathcal B(q,r)
 ={}&\beta^2K\alpha(q)\bigl(q-\Gamma(q)\bigr)\\
 &+2\beta^2\int_{s_1}^{s_2}\alpha(s)k(s)
   \int_{s_1}^{s}k(t)
     \left\{\beta^2\E\bigl[\chi(q,X_q^x)V_t\bigr]-1\right\}
           \dd t\dd s.
 \end{split}
\end{equation}
Decrease $\tau$ further until $\beta^2K\tau\leq\gamma W$. The first
term in~\eqref{static-eq:locking-bracket} is then at most $\gamma W$, and
\eqref{static-eq:locking-negative-factor} bounds the double integral
by $-2\gamma W$. Hence
\[
 \mathcal B(q,r)\leq\gamma W-2\gamma W=-\gamma W.
\]
The cubic remainder in~\eqref{static-eq:locking-three-pressure-expansion}
is uniform over these $q,r$. Choose $\delta_0>0$ small enough to
absorb it. For some $c_0>0$ we have obtained
\begin{equation}
\label{static-eq:locking-three-gap}
 F_{3,n}(D)\leq3P_\beta-c_0\delta^2
\end{equation}
whenever the constraint is nonempty,
$0\leq r\leq r_0$, $|\delta|\leq\delta_0$, $q\in[h,1]$, and
$|q-\Gamma(q)|\leq\tau$. In particular, this estimate places no
restriction $r\geq|\delta|$ on the two small signed overlaps.

\subsection{From the pressure gap to locking}

\begin{proof}[Proof of Lemma~\ref{lem:fixed-locking}]
Choose $b\le\min(r_0,\delta_0)$. For a triple in
\eqref{eq:fixed-locking}, set $q=R(x,y)\ge h$ and reverse $v$ if
necessary to make
\[
 r=\frac{R(v,x)+R(v,y)}2\ge0,
 \qquad \delta=\frac{R(v,x)-R(v,y)}2.
\]
This reversal preserves the Gibbs weight, the large overlap $q$, and
$|\delta|$. The event gives $r,|\delta|\le b$ and
$|\delta|>n^{-e_*}$. In particular, we do not need the two small
overlaps to have the same sign.

Lemma~\ref{static-lem:pair-bound} bounds pairs with
$|q-\Gamma(q)|>\tau$ by $C\exp(-cn)$ in averaged probability.
Adding an independent Gibbs draw does not change that bound.
For every remaining feasible triple,
\eqref{static-eq:locking-three-gap} gives a pressure deficit
$\eta=c_0\delta^2\ge c_0n^{-2e_*}$. Since
$2e_*=1/50<1/24$, this dominates the finite-size error in
Lemma~\ref{lem:pressure-prob}, uniformly over the constraints.
Their averaged probabilities are bounded by
\[
 \exp(-c n^{49/50})+\exp(-c n^{24/25}).
\]
There are at most $(n+1)^3$ overlap triples and two choices of the
reversal. The resulting bound is at most
$C(n+1)^3\exp(-c n^{24/25})$, which is smaller than
$\exp(-n^{9/10})$ for all sufficiently large $n$.
\end{proof}

\begin{proposition}[Absolute-overlap locking]
\label{static-prop:locking}
For every fixed $h\in(0,1)$ there is $b>0$ such that, for all
sufficiently large $n$,
\begin{equation}\label{static-eq:locking-probability}
 \E\pi^{\otimes3}\!\left\{
 \begin{gathered}
 |R(x,y)|\ge h,\qquad |R(v,x)|,|R(v,y)|\le b,\\
 \bigl||R(v,x)|-|R(v,y)|\bigr|>2n^{-1/100}
 \end{gathered}\right\}\le \exp(-n^{9/10}).
\end{equation}
The same conclusion holds after decreasing $b$.
\end{proposition}
\begin{proof}
Take $b$ from Lemma~\ref{lem:fixed-locking}. Reverse $y$ if needed to
make $R(x,y)\ge h$. This preserves the Gibbs weight and all absolute
overlaps in the event. The inequality
\[
 \bigl||R(v,x)|-|R(v,y)|\bigr|
       \le |R(v,x)-R(v,y)|
\]
then puts the transformed triple in the signed event. The two possible
reversals multiply its probability bound by at most two. The stronger
bound $C(n+1)^3\exp(-c n^{24/25})$ obtained in the preceding proof
absorbs this factor and gives the claimed exponent.
\end{proof}
